\chapter*{Résumé}
\addcontentsline{toc}{chapter}{Résumé}

Ce mémoire présente la conception de la partie du projet consacrée à la
détection de signaux techniques et à leur restitution dans une interface de
supervision destinée au marché marocain. L'enjeu principal est de produire des
signaux exploitables sans se limiter à une logique de performance apparente,
mais en privilégiant la robustesse hors échantillon, la stabilité temporelle et
l'explicabilité des résultats.

Le document se concentre sur trois éléments complémentaires. Le premier est la
méthodologie de construction des signaux, fondée sur une chaîne de traitement en
plusieurs étapes : génération de variantes candidates, évaluation hors
échantillon, filtrage de robustesse, réduction de redondance puis agrégation au
niveau des familles. Le deuxième est la version initiale de la page
\texttt{/signals}, qui permet une lecture progressive des résultats depuis une
vue synthétique jusqu'au détail d'une famille donnée. Le troisième est le
tableau de bord, qui transforme ces sorties en un dispositif de hiérarchisation
des opportunités par horizon et par titre.

La contribution principale de cette partie du projet réside dans la mise en
cohérence d'une chaîne complète allant des données de marché à une restitution
opérationnelle compréhensible par l'utilisateur final.

\textbf{Mots-clés :} signaux techniques, robustesse, hors échantillon, tableau
de bord, marché marocain.

\cleardoublepage

